% !PRINTED 534
\weberpage{534}

%<<CONT
which \gesperrt{Abel} taught how to solve in general, and which we shall therefore, after him, call \gesperrt{Abelian equations}\footnote{\gesperrt{Abel}, Mémoire sur une classe d'équations résolubles algébriquement. Crelle's Journal f. Mathematik, Bd.~4, 1829. Oeuvres complètes, nouvelle édition, 1881, Bd.~1, S.~418.}.

\gesperrt{An equation of degree $m$, $F(x) = 0$, with the roots $\alpha, \alpha_1, \alpha_2 \ldots \alpha_{m-1}$, is called an Abelian equation if every root is rationally expressible in terms of one of them, $\alpha$, and if, when}
\weq{2}{\alpha_1 = \Theta_1(\alpha), \quad \alpha_2 = \Theta_2(\alpha) \ldots \alpha_{m-1} = \Theta_{m-1}(\alpha)}
\gesperrt{are these rational expressions, the condition}
\weq{3}{\Theta_h\Theta_k(\alpha) = \Theta_k\Theta_h(\alpha)}
\gesperrt{holds for every pair of these functions.}

Here the symbol $\Theta_h\Theta_k(\alpha)$ means that the function $\Theta_h(x)$ is to be formed for the argument $x = \Theta_k(\alpha)$. This whole definition refers, of course, to a definite field $\Omega$ assumed as rational. To these equations certainly belong, when $\Omega$ is the field of rational numbers, the equations by which the roots of unity are determined. For every $m$th root of unity is, if $r$ is a primitive one among them, contained in the form $\Theta_h(r) = r^h$, and we have
\[
\Theta_h\Theta_k(r) = \Theta_k\Theta_h(r) = r^{hk}.
\]
Likewise, when the field $\Omega$ contains the $m$th roots of unity $r$, the \gesperrt{pure equation}
\[
x^m - a = 0,
\]
in which $a$ belongs to the field $\Omega$, belongs to the \gesperrt{Abelian equations}. For if $\alpha$ is one of its roots, then all the roots are contained in the form
\[
\alpha_h = r^h\alpha
\]
we have therefore $\Theta_h(x) = r^hx$, and hence
\[
\Theta_h\Theta_k(x) = \Theta_k\Theta_h(x) = r^{h+k}x.
\]
If the function $F(x)$ is not irreducible, then it has a definite irreducible factor $\varphi(x)$ which has the root $\alpha$, and among the roots of $\varphi(x) = 0$ the relations expressed by (2) and (3) likewise hold. Hence $\varphi(x) = 0$ is, by \S.\,145, a \gesperrt{Galois} resolvent of
%CONT>>
